\chapter{ویرایشگر متنی vi(m)}

در میان اهالی نیویورک، لطیفه‌ی قدیمی‌ای رایج است: مسافری از رهگذری می‌پرسد: «ببخشید، چگونه می‌توان به تالار کارنگی (محل اجرای موسیقی کلاسیک) رسید؟» و رهگذر پاسخ می‌دهد: «با تمرین، تمرین و باز هم تمرین!». تسلط بر خط فرمان لینوکس نیز شباهت زیادی به مهارت در نواختن پیانو دارد؛ توانایی‌ای که در یک بعدازظهر به دست نمی‌آید، بلکه حاصل سال‌ها ممارست و تجربه است. در این فصل به معرفی ویرایشگر متن \cmd{vi} (تلفظ: وی‌آی) می‌پردازیم؛ ابزاری بنیادین و کلاسیک در فرهنگ یونیکس که با وجود شهرت دوگانه‌اش به‌سبب رابط کاربری خاص و متفاوت، در دستان کاربری حرفه‌ای، اجرایی هنرمندانه بر صفحه‌کلید را تداعی می‌کند. مقصود ما در این فصل، نه دستیابی به درجهٔ استادی، بلکه توانایی اجرای روانِ دستورات اولیه و ضروری این ویرایشگر است.

\section{ضرورت یادگیری vi}

در دورانی که ویرایشگرهای گرافیکی (\en{GUI}) قدرتمند و ابزارهای متنی ساده‌ای مانند \cmd{nano} به‌راحتی در دسترس‌اند، این پرسش مطرح می‌شود که چرا همچنان باید وقت خود را صرف یادگیری \cmd{vi} کنیم؟ پاسخ به این سؤال را می‌توان در سه دلیل اصلی خلاصه کرد:

\begin{itemize}
  \item \textbf{دسترس‌پذیری همه‌جانبه:} \cmd{vi} تقریباً در تمامی توزیع‌های لینوکس و سیستم‌های شبه‌یونیکس به‌صورت پیش‌فرض نصب است. این ویژگی در سناریوهای بحرانی، مانند مدیریت سرورهای راه دور فاقد محیط گرافیکی یا تعمیر سیستم‌هایی که رابط بصری آن‌ها آسیب دیده، یک مزیت حیاتی محسوب می‌شود. در حالی که \cmd{nano} محبوبیت فزاینده‌ای یافته است، اما هنوز جزو الزامات استاندارد \en{POSIX} نیست؛ در مقابل، اگرچه \cmd{vi} به‌عنوان یک «گزینه» (و نه یک الزام) در این استاندارد تعریف شده، اما در عمل، تقریباً بر روی تمامی سیستم‌های سازگار با یونیکس موجود است و به‌عنوان یک ویرایشگر استاندارد و همیشگی در این محیط‌ها شناخته می‌شود.

  \item \textbf{سبکی و سرعت اجرا:} در بسیاری از موارد، فراخوانی \cmd{vi} بسیار سریع‌تر از جست‌وجو در منوهای گرافیکی و انتظار برای بارگذاری برنامه‌های حجیم است. فراتر از آن، \cmd{vi} برای بیشترین سرعت در تایپ طراحی شده است. همان‌طور که در ادامه خواهیم دید، یک کاربر ماهر می‌تواند بدون جدا کردن انگشتان از ردیف اصلی صفحه‌کلید و بدون نیاز به ماوس، عملیات ویرایشی پیچیده را با سرعتی چشمگیر انجام دهد.

  \item \textbf{اعتبار در میان جامعه کاربران:} طبیعتاً هیچ‌کس دوست ندارد در میان سایر متخصصان لینوکس و یونیکس، به‌عنوان کاربری ناآشنا با ابزارهای اصیل این حوزه شناخته شود! البته، شاید همان دو دلیل نخست برای متقاعد کردن هر کاربر حرفه‌ای کافی باشد!
\end{itemize}

\section{پیشینه و سیر تحول vi}

نخستین نسخه از \cmd{vi} در سال ۱۹۷۶ توسط بیل جوی (\en{Bill Joy})، دانشجوی وقت دانشگاه برکلی و یکی از بنیان‌گذاران آتی شرکت \en{Sun Microsystems}، توسعه یافت. نام‌گذاری این ویرایشگر برگرفته از واژه‌ی \en{Visual} است؛ چراکه برخلاف نسل‌های پیشین خود، برای کار بر روی ترمینال‌های تصویری مجهز به مکان‌نمای متحرک طراحی شده بود. پیش از ظهور \cmd{vi}، تعامل با متن از طریق «ویرایشگرهای خطی» (\en{Line Editors}) صورت می‌گرفت که در هر لحظه تنها بر روی یک سطر متمرکز بودند. در آن سیستم‌ها، کاربر ناچار بود ابتدا آدرس سطر مورد نظر را فراخوانی کرده و سپس دستور تغییر، حذف یا درج را صادر کند. با جایگزینی ترمینال‌های تصویری به‌جای دستگاه‌های تله‌تایپ (\en{Teletype})، امکان ویرایش بلادرنگ و بصری متن فراهم گردید.

جالب است بدانید که \cmd{vi} در بطن خود یک ویرایشگر خطی قدرتمند به نام \cmd{ex} را نیز در بر دارد و کاربران حرفه‌ای در حین کار با \cmd{vi}، به‌طور مداوم از دستورات \cmd{ex} برای انجام عملیات دسته‌جمعی و پیچیده استفاده می‌کنند.

امروزه در اکثر توزیع‌های لینوکس، به‌جای نسخه سنتی و اولیه، از نسخه‌ای ارتقایافته به نام \cmd{vim} (مخفف \en{vi improved}) استفاده می‌شود که توسط برام مولنار (\en{Bram Moolenaar}) خلق شده است. \cmd{vim} نسبت به نسخه کلاسیک یونیکس پیشرفت‌های شگرفی داشته و معمولاً در سیستم‌های لینوکسی از طریق یک پیوند نمادین (\en{Symbolic Link}) یا نام مستعار (\en{Alias})، با همان فرمان ساده‌ی \cmd{vi} فراخوانی می‌شود.

در ادامه‌ی این مباحث، فرض ما بر این است که با اجرای دستور \cmd{vi}، در حقیقت از توانمندی‌های نسخه پیشرفته یعنی \cmd{vim} بهره‌مند می‌شوید.

\section{فراخوانی و خروج از محیط vi}

برای اجرای ویرایشگر \cmd{vi} کافی است نام آن را در خط فرمان وارد نمایید:

\begin{latin}
\begin{lstlisting}
[me@linuxbox ~]$ vi
\end{lstlisting}
\end{latin}

پس از اجرای دستور، با رابط کاربری متنی مشابه تصویر زیر مواجه خواهید شد. نویسه‌های تیلدا (\en{Tilde} یا \cmd{\textasciitilde}) در ابتدای سطرها نشان‌دهنده آن است که این سطور خالی هستند و هنوز متنی در آن‌ها درج نشده است.

\begin{latin}
\begin{lstlisting}
~
~
~                           VIM - Vi Improved
~                            version 8.0.707
~                        by Bram Moolenaar et al.
~               Vim is open source and freely distributable
~
~                        Sponsor Vim development!
~             type :help sponsor<Enter>    for information
~
~             type :q<Enter>                 to exit
~             type :help<Enter> or <F1>      for on-line help
~             type :help version8<Enter>     for version info
~             Running in Vi compatible mode
~             type :set nocp<Enter>          for Vim defaults
~             type :help cp-default<Enter>   for info on this
~
~
\end{lstlisting}
\end{latin}

همان‌طور که در کار با ویرایشگر \cmd{nano} تجربه کردیم، نخستین گام حیاتی در مواجهه با هر ابزار جدید، فراگیری روش صحیح خروج از آن است. برای خروج از \cmd{vi}، دستور زیر را تایپ کنید (دقت داشته باشید که نویسه‌ی «دو نقطه» پیشوند الزامی این دستور است):

\begin{latin}
\begin{lstlisting}
:q
\end{lstlisting}
\end{latin}

پس از وارد کردن این دستور و فشردن کلید \en{Enter}، برنامه بسته شده و مجدداً اعلانِ پوسته (\en{Shell Prompt}) ظاهر می‌گردد. چنانچه فرآیند خروج با موفقیت انجام نشد، این امر معمولاً به معنای وجود تغییرات ذخیره‌نشده در فایل است. در چنین شرایطی، اگر تمایلی به حفظ تغییرات ندارید، می‌توانید با افزودن علامت تعجب (\cmd{!}) به انتهای دستور، خروج اجباری (\en{Force Quit}) را به برنامه تحمیل کنید:

\begin{latin}
\begin{lstlisting}
:q!
\end{lstlisting}
\end{latin}

\begin{note}
\textbf{نکته:} چنانچه در محیط \cmd{vi} با وضعیت غیرمنتظره‌ای مواجه شدید یا اصطلاحاً در میان حالت‌های مختلف آن «سرگردان» شدید، با دو بار فشردن کلید \en{Esc} می‌توانید اطمینان حاصل کنید که برنامه به حالت فرمان (\en{Command Mode}) بازگشته است. این عمل وضعیت ویرایشگر را به حالت پیش‌فرض بازگردانده و کنترل محیط را مجدداً در اختیار شما قرار می‌دهد.
\end{note}

\section{حالت سازگاری (Compatibility Mode) در vim}

در رابط کاربری اولیه‌ی برنامه که پیش‌تر بررسی کردیم، عبارت \en{``Running in Vi compatible mode''} مشاهده می‌شود. این پیام به این معناست که \cmd{vim} در وضعیتی اجرا شده است که رفتار نسخه‌ی کلاسیک و قدیمی \cmd{vi} را شبیه‌سازی می‌کند و بسیاری از قابلیت‌های بهبودیافته‌ی خود را غیرفعال کرده است.

به منظور بهره‌مندی از حداکثر پتانسیل \cmd{vim} در تمرین‌های این فصل، توصیه می‌شود آن را در حالت بهبودیافته اجرا کنید. برای این منظور، می‌توانید یکی از راهکارهای زیر را به کار بگیرید:

نخست، تلاش کنید به جای دستور \cmd{vi}، مستقیماً از دستور \cmd{vim} استفاده کنید. در صورت موفقیت‌آمیز بودن این روش، می‌توانید با افزودن سطر زیر به فایل پیکربندی \file{.bashrc}، یک نام مستعار دائمی ایجاد کنید:

\begin{latin}
\begin{lstlisting}
alias vi='vim'
\end{lstlisting}
\end{latin}

همچنین می‌توانید با اجرای دستور زیر در ترمینال، گزینه‌ی غیرفعال‌سازی حالت سازگاری (\cmd{set nocp}) را مستقیماً به فایل تنظیمات اختصاصی این ویرایشگر (\file{.vimrc}) الحاق کنید:

\begin{latin}
\begin{lstlisting}
echo "set nocp" >> ~/.vimrc
\end{lstlisting}
\end{latin}

شایان ذکر است که توزیع‌های مختلف لینوکس، ابزار \cmd{vim} را در قالب بسته‌های نرم‌افزاری متفاوتی عرضه می‌کنند. برخی از توزیع‌ها به صورت پیش‌فرض نسخه‌های سبک و حداقلی مانند \cmd{vim-tiny} را نصب می‌کنند که تنها از مجموعه‌ی محدودی از قابلیت‌ها پشتیبانی می‌کند.

چنانچه در طول انجام تمرین‌ها متوجه شدید برخی از ویژگی‌های معرفی‌شده در سیستم شما در دسترس نیستند، پیشنهاد می‌شود نسخه‌ی کامل و استاندارد \cmd{vim} را بر روی توزیع خود نصب نمایید.

\section{مدیریت حالت‌های عملیاتی (Modal Editing)}

بیایید مجدداً ویرایشگر \cmd{vi} را فراخوانی کنیم؛ این بار نام فایلی را وارد می‌کنیم که هنوز وجود ندارد. این یکی از روش‌های رایج برای ایجاد یک فایل جدید در محیط ترمینال است:

\begin{latin}
\begin{lstlisting}
[me@linuxbox ~]$ rm -f foo.txt
[me@linuxbox ~]$ vi foo.txt
\end{lstlisting}
\end{latin}

در صورت اجرای صحیح، با محیطی مشابه زیر روبه‌رو خواهید شد:

\begin{latin}
\begin{lstlisting}
~
~
~
~
~
~
~
~
~
~
~
~
~
~
~
~
~
~

"foo.txt" [New File]
\end{lstlisting}
\end{latin}

نویسه‌های تیلدا (\cmd{\textasciitilde}) در ستون سمت چپ، نشان‌دهنده سطرهای تهی هستند و تأیید می‌کنند که فایل در حال حاضر فاقد محتواست. در این مرحله، اکیداً توصیه می‌شود پیش از درک ساختار برنامه، هیچ کلیدی را فشار ندهید.

پس از یادگیری روش خروج، حیاتی‌ترین مفهومی که باید درک کنید «حالت‌مند» بودن (\en{Modal}) این ویرایشگر است. برخلاف ویرایشگرهای متداول، \cmd{vi} هنگام اجرا به‌صورت پیش‌فرض در حالت عادی (\en{Normal Mode}) قرار می‌گیرد. در این وضعیت، تقریباً هر کلیدی که روی صفحه‌کلید فشار می‌دهید، به جای درج متن، به‌عنوان یک «فرمان» (\en{Command}) تفسیر می‌شود. بنابراین، اگر بدون تغییر حالت شروع به تایپ کنید، \cmd{vi} رفتارهای غیرمنتظره‌ای از خود نشان داده و ممکن است چیدمان متن یا تنظیمات فایل شما را به کلی دگرگون کند.

\subsection{ورود به حالت درج (Insert Mode)}

به منظور افزودن محتوا به فایل، نخست باید از حالت عادی خارج شده و وارد حالت درج (\en{Insert Mode}) شوید. برای این انتقال، کافی است کلید \cmd{i} را بر روی صفحه‌کلید فشار دهید. چنانچه \cmd{vim} در حالت بهبودیافته (و نه حالت سازگار با نسخه کلاسیک) در حال اجرا باشد، نشانگر زیر را در پایین‌ترین سطر صفحه مشاهده خواهید کرد:

\begin{latin}
\begin{lstlisting}
-- INSERT --
\end{lstlisting}
\end{latin}

اکنون محیط برای تایپ متن آماده است. می‌توانید جمله نمونه زیر را وارد کنید:

\begin{latin}
\begin{lstlisting}
The quick brown fox jumps over the lazy dog.
\end{lstlisting}
\end{latin}

پس از اتمام تایپ، برای خروج از این وضعیت و بازگشت به حالت عادی (\en{Normal Mode}) جهت اجرای فرمان‌ها، کلید \en{Esc} را فشار دهید. با این کار، عبارت \cmd{-- INSERT --} از پایین صفحه حذف شده و مجدداً در حالت فرمان قرار خواهید گرفت.

\subsection{ذخیره‌سازی تغییرات}

برای ثبت تغییرات اعمال‌شده در فایل بر روی دیسک، ابتدا باید وارد حالت خط فرمان (\en{Command-line Mode}) شوید. این وضعیت با فشردن کلید \cmd{:} (دو نقطه) در حالت عادی (\en{Normal Mode}) فعال می‌شود. بلافاصله پس از فشردن این کلید، نویسه‌ی \cmd{:} در پایین‌ترین سطر صفحه ظاهر شده و منتظر دریافت دستور می‌ماند:

\begin{latin}
\begin{lstlisting}
:
\end{lstlisting}
\end{latin}

به منظور نوشتن (ذخیره‌سازی) محتوا بر روی فایل، پس از علامت دو نقطه، حرف \cmd{w} (مخفف \en{Write}) را تایپ کرده و کلید \en{Enter} را فشار دهید:

\begin{latin}
\begin{lstlisting}
:w
\end{lstlisting}
\end{latin}

با انجام این عمل، فایل بر روی دیسک ذخیره شده و گزارشی از وضعیت آن در پایین صفحه نمایش داده می‌شود که نشان‌دهنده‌ی نام فایل، تعداد سطور و تعداد نویسه‌های ذخیره‌شده است:

\begin{latin}
\begin{lstlisting}
"foo.txt" [New] 1L, 45C written
\end{lstlisting}
\end{latin}

\begin{note}
\textbf{نکته‌ای در باب اصطلاح‌شناسی:} لازم است بدانید که میان مستندات \cmd{vim} و نسخه‌ی کلاسیک \cmd{vi} در نام‌گذاری حالت‌ها تفاوت‌هایی وجود دارد. \cmd{vim} سه حالت اصلی خود را با عناوین زیر می‌شناسد:
\begin{enumerate}
  \item \en{Normal}
  \item \en{Insert}
  \item \en{Command}
\end{enumerate}
در مقابل، منابع مربوط به نسخه‌ی اصلی \cmd{vi} (و بسیاری از مراجع قدیمی‌تر) همین سه حالت را عمدتاً به این صورت خطاب می‌کنند:
\begin{enumerate}
  \item \en{Command}
  \item \en{Insert}
  \item \en{ex}
\end{enumerate}
آگاهی از این تفاوتِ نام‌گذاری برای جلوگیری از سردرگمی هنگام مطالعه‌ی مراجع مختلف ضروری است؛ به‌ویژه آنچه در \cmd{vim} با عنوان \en{Normal} شناخته می‌شود، در منابع کلاسیک اغلب \en{Command} نامیده می‌شود.

همچنین در مورد حالت سوم باید توجه داشت که در \cmd{vim}، حالت \en{Command-line} (که با فشردن \cmd{:} فعال می‌شود) و حالت \en{Ex} (که با فشردن \cmd{Q} فعال می‌شود) دو حالت مجزا و مستقل از یکدیگرند؛ اما از آنجا که در \cmd{vi} کلاسیک، فشردن \cmd{:} کاربر را مستقیماً به محیط خط‌فرمانِ ویرایشگر \cmd{ex} می‌برد، بسیاری از منابع قدیمی این حالت را با عنوان \en{ex mode} می‌شناسند. شایان ذکر است که برخی مراجع کلاسیک، به‌جای عبارت \en{ex mode}، از اصطلاح \en{command-line mode} نیز برای توصیف همین حالت استفاده کرده‌اند و از این‌رو نام‌گذاری در منابع مختلف کاملاً یکدست نیست.
\end{note}

\section{پیمایش و جابه‌جایی مکان‌نما}

در حالت عادی (\en{Normal Mode})، ویرایشگر \cmd{vi} طیف گسترده‌ای از دستورات را برای جابه‌جایی مکان‌نما در اختیار کاربر قرار می‌دهد که بسیاری از آن‌ها با کلیدهای پیمایش در ابزار \cmd{less} مشابهت دارند. جدول زیر، زیرمجموعه‌ای از پرکاربردترین دستورات حرکتی را نمایش می‌دهد:

\begin{table}[htbp]
\centering
\caption{کلیدهای حرکت مکان‌نما در حالت عادی \cmd{vi}}
\label{tab:vi-movement}
\begin{tabularx}{\textwidth}{l X}
\toprule
\textbf{کلید} & \textbf{عملکرد فنی در حالت عادی (\en{Normal Mode})} \\
\midrule
\cmd{l} یا $\rightarrow$ & انتقال مکان‌نما به اندازه‌ی یک نویسه به سمت راست \\
\addlinespace
\cmd{h} یا $\leftarrow$ & انتقال مکان‌نما به اندازه‌ی یک نویسه به سمت چپ \\
\addlinespace
\cmd{j} یا $\downarrow$ & انتقال مکان‌نما به سطر پایین \\
\addlinespace
\cmd{k} یا $\uparrow$ & انتقال مکان‌نما به سطر بالا \\
\addlinespace
\cmd{0} (صفر) & انتقال به ابتدای مطلق سطر جاری \\
\addlinespace
\cmd{\^{}} & پرش به نخستین نویسه‌ی غیرفضای‌خالی (\en{Non-whitespace}) در سطر جاری \\
\addlinespace
\cmd{\$} & پرش به انتهای سطر جاری \\
\addlinespace
\cmd{w} & انتقال به ابتدای کلمه یا نویسه‌ی نشانه‌گذاری (علائم نگارشی - \en{Punctuation}) بعدی \\
\addlinespace
\cmd{W} & انتقال به ابتدای کلمه بعدی (با نادیده گرفتن نویسه‌های نشانه‌گذاری) \\
\addlinespace
\cmd{b} & انتقال به ابتدای کلمه یا نویسه‌ی نشانه‌گذاریِ قبلی \\
\addlinespace
\cmd{B} & انتقال به ابتدای کلمه قبلی (با نادیده گرفتن نویسه‌های نشانه‌گذاری) \\
\addlinespace
\cmd{Ctrl+f} یا \en{Page Down} & پیمایش یک صفحه کامل به سمت پایین (\en{Forward}) \\
\addlinespace
\cmd{Ctrl+b} یا \en{Page Up} & پیمایش یک صفحه کامل به سمت بالا (\en{Backward}) \\
\addlinespace
\cmd{numberG} & انتقال مستقیم به سطر شماره \file{number} (مثلاً \cmd{1G} برای پرش به سطر اول) \\
\addlinespace
\cmd{G} & انتقال مستقیم به آخرین سطر فایل \\
\bottomrule
\end{tabularx}
\end{table}

\begin{note}
\textbf{نکته:} در علوم رایانه، به مجموعه‌ای از نویسه‌ها که در متن قابل مشاهده نیستند اما فضایی را اشغال می‌کنند، نویسه‌های فضای خالی (\en{Whitespace Characters}) گفته می‌شود. این نویسه‌ها برای ایجاد فاصلهٔ افقی یا عمودی میان دیگر نویسه‌ها به کار می‌روند و رایج‌ترین آن‌ها عبارت‌اند از: فاصله (\cmd{' '})، تب (\cmd{\textbackslash t})، خط جدید (\cmd{\textbackslash n}) و بازگشت به ابتدای سطر (\cmd{\textbackslash r}). در مقابل، نویسه‌های غیرفضای‌خالی (\en{Non-whitespace Characters}) به هر نویسه‌ای گفته می‌شود که در دستهٔ فوق قرار نگیرد؛ از جمله تمام حروف (مانند \cmd{a}، ب)، اعداد (مانند \cmd{1}، ۲) و نمادها (مانند \cmd{@}، \cmd{\#}، \cmd{.}). به بیان ساده، هر نویسه‌ای که قابل مشاهده باشد و فضای خالی ایجاد نکند، یک نویسه‌ی غیرفضای‌خالی محسوب می‌شود.

\textbf{نکته:} در ویرایشگر \cmd{vi} و \cmd{vim}، نویسه‌های نشانه‌گذاری (\en{Punctuation}) مانند کاما (\cmd{,})، نقطه (\cmd{.})، سمی‌کالن (\cmd{;}) و علامت سؤال (\cmd{?}) به‌عنوان «کلمه» (\en{word}) در نظر گرفته می‌شوند. به همین دلیل، دستور \cmd{w} هنگام حرکت، روی این نویسه‌ها نیز توقف می‌کند و آن‌ها را به‌عنوان یک واحد مجزا برای حرکت در نظر می‌گیرد. در مقابل، نسخهٔ بزرگ آن یعنی \cmd{W} (و همچنین \cmd{B} و \cmd{E}) این نویسه‌ها را نادیده گرفته و تنها بر اساس نویسه‌های فضای خالی (\en{whitespace characters}) مانند فاصله (\cmd{' '})، تب (\cmd{\textbackslash t}) و خط جدید (\cmd{\textbackslash n}) حرکت می‌کند.
\end{note}

دلیل انتخاب کلیدهای \cmd{h}، \cmd{j}، \cmd{k} و \cmd{l} برای جابه‌جایی مکان‌نما در ویرایشگر \cmd{vi}، نخست به محدودیت‌های سخت‌افزاری ترمینال‌های اولیه بازمی‌گردد؛ چرا که بسیاری از آن‌ها فاقد کلیدهای جهت‌نمای اختصاصی بودند. افزون بر این، این چیدمان به کاربران حرفه‌ای امکان می‌دهد تا بدون خارج کردن انگشتان از ردیف اصلی صفحه‌کلید (\en{Home Row})، مکان‌نما را با سرعت و دقت بیشتری کنترل کنند.

بسیاری از دستورات \cmd{vi} قابلیت پذیرش یک «پیشوند عددی» را دارند. با استفاده از این ویژگی، می‌توان تعداد دفعات تکرار یک دستور را پیش از اجرای آن تعیین کرد. به عنوان نمونه، وارد کردن دستور \cmd{5j} موجب می‌شود تا مکان‌نما به‌صورت آنی پنج سطر به سمت پایین حرکت کند.

\section{عملیات ویرایشی مقدماتی}

اغلب فرآیندهای ویرایشی بر چند عملیات کلیدی استوارند: درج محتوا، حذف نویسه‌ها و جابه‌جایی متن از طریق برش و چسباندن (\en{Cut and Paste}). ویرایشگر \cmd{vi} تمامی این قابلیت‌ها را با منطق منحصربه‌فرد خود پیاده‌سازی می‌کند. علاوه بر این، \cmd{vi} از قابلیت «واگرد» (\en{Undo}) نیز پشتیبانی می‌کند؛ به گونه‌ای که با فشردن کلید \cmd{u} در حالت عادی (\en{Normal Mode})، آخرین تغییر اعمال شده لغو و فایل به وضعیت پیشین بازمی‌گردد. بهره‌گیری از این ویژگی در هنگام اجرای دستورات ویرایشی، نقشی حیاتی در مدیریت خطاها ایفا می‌کند.

\subsection{الحاق متن}

ویرایشگر \cmd{vi} روش‌های متعددی را برای ورود به «حالت درج» (\en{Insert Mode}) در اختیار کاربر قرار می‌دهد. پیش‌تر از دستور \cmd{i} برای درج متن در موقعیت فعلی مکان‌نما استفاده کردیم.

بیایید مجدداً فایل \file{foo.txt} را بررسی کنیم:

\begin{latin}
\begin{lstlisting}
The quick brown fox jumps over the lazy dog.
\end{lstlisting}
\end{latin}

چنانچه بخواهیم متنی را به انتهای این جمله اضافه کنیم، متوجه خواهیم شد که دستور \cmd{i} به‌تنهایی پاسخگو نیست؛ زیرا در حالت عادی، مکان‌نما اجازه عبور از آخرین نویسه‌ی سطر را ندارد. برای رفع این محدودیت، \cmd{vi} دستور \cmd{a} (مخفف \en{append}) را ارائه می‌دهد. با انتقال مکان‌نما به انتهای سطر و فشردن کلید \cmd{a}، مکان‌نما یک واحد به جلو حرکت کرده و برنامه وارد حالت درج می‌شود. اکنون می‌توانیم متن مورد نظر را الحاق کنیم:

\begin{latin}
\begin{lstlisting}
The quick brown fox jumps over the lazy dog. It was cool.
\end{lstlisting}
\end{latin}

به خاطر داشته باشید که برای خروج از حالت درج و بازگشت به حالت عادی، فشردن کلید \en{Esc} الزامی است.

به دلیل نیاز مکرر به الحاق متن در انتهای سطر، \cmd{vi} میان‌بر کارآمد \cmd{A} را ارائه می‌دهد. این دستور، مکان‌نما را مستقیماً به انتهای سطر جاری منتقل کرده و بلافاصله وضعیت را به «حالت درج» (\en{Insert Mode}) تغییر می‌دهد.

بیایید این قابلیت را با افزودن سطرهای جدید به فایل تمرین کنیم: نخست با فشردن کلید \cmd{0} (صفر)، مکان‌نما را به ابتدای سطر بازگردانید؛ سپس با تایپ \cmd{A}، متن زیر را وارد نمایید:

\begin{latin}
\begin{lstlisting}
The quick brown fox jumps over the lazy dog. It was cool.
Line 2
Line 3
Line 4
Line 5
\end{lstlisting}
\end{latin}

پس از اتمام تایپ، با فشردن کلید \en{Esc} به حالت عادی بازگردید. همان‌طور که مشاهده کردید، دستور \cmd{A} به مراتب کارآمدتر از \cmd{a} عمل می‌کند، زیرا فرآیند انتقال به انتهای سطر و ورود به حالت درج را در یک گام ادغام می‌نماید.

\subsection{درج سطر جدید}

یکی دیگر از شیوه‌های ورود به «حالت درج» (\en{Insert Mode})، استفاده از قابلیت «باز کردن سطر» (\en{Open}) است. این عملیات یک سطر خالی میان دو سطر موجود ایجاد کرده و محیط را آماده‌ی دریافت متن می‌کند. این قابلیت در دو حالت کلی مطابق جدول زیر عمل می‌کند:

\begin{table}[htbp]
\centering
\caption{کلیدهای درج سطر جدید}
\label{tab:vi-new-line}
\begin{tabularx}{\textwidth}{c X}
\toprule
\textbf{دستور} & \textbf{عملکرد} \\
\midrule
\cmd{o} & ایجاد سطر جدید در پایین سطر فعلی \\
\addlinespace
\cmd{O} & ایجاد سطر جدید در بالای سطر فعلی \\
\bottomrule
\end{tabularx}
\end{table}

برای درک عملی این موضوع، مکان‌نما را بر روی عبارت «\en{Line 3}» قرار داده و کلید \cmd{o} (حرف کوچک) را فشار دهید:

\begin{latin}
\begin{lstlisting}
The quick brown fox jumps over the lazy dog. It was cool.
Line 2
Line 3

Line 4
Line 5
\end{lstlisting}
\end{latin}

مشاهده می‌کنید که یک سطر تهی در زیر سطر سوم ایجاد شده و ویرایشگر مستقیماً وارد حالت درج شده است. جهت اتمام عملیات، کلید \en{Esc} را بفشارید. برای بازگرداندن فایل به وضعیت پیشین (\en{Undo})، کلید \cmd{u} را بزنید.

حال، بار دیگر مکان‌نما را روی «\en{Line 3}» قرار داده و این بار کلید \cmd{O} (حرف بزرگ) را فشار دهید تا سطری در بالای موقعیت فعلی باز شود:

\begin{latin}
\begin{lstlisting}
The quick brown fox jumps over the lazy dog. It was cool.
Line 2

Line 3
Line 4
Line 5
\end{lstlisting}
\end{latin}

برای خروج از حالت درج، کلید \en{Esc} و برای لغو تغییرات اعمال شده، کلید \cmd{u} را فشار دهید.

\subsection{حذف محتوا}

ویرایشگر \cmd{vi} راهکارهای متنوعی را برای حذف متن ارائه می‌دهد که زیربنای همگی آن‌ها بر دو کلید اصلی استوار است: دستور \cmd{x} برای حذف نویسه (\en{Character}) در محل استقرار مکان‌نما به کار می‌رود. این دستور می‌تواند با یک پیشوند عددی ترکیب شود تا تعداد مشخصی از نویسه‌ها را به‌طور هم‌زمان حذف کند. دستور \cmd{d} کاربرد گسترده‌تری دارد؛ این فرمان علاوه بر قابلیت پذیرش عدد، باید با یک «دستور حرکتی» (مانند \cmd{w}، \cmd{\$}، یا \cmd{G}) ترکیب شود تا بازه‌ی عملیاتی (\en{Scope}) حذف را تعیین نماید.

\begin{table}[htbp]
\centering
\caption{دستورات حذف متن}
\label{tab:vi-delete}
\begin{tabularx}{\textwidth}{l X}
\toprule
\textbf{دستور} & \textbf{عملکرد فنی} \\
\midrule
\cmd{x} & حذف نویسه‌ی جاری \\
\addlinespace
\cmd{3x} & حذف نویسه‌ی جاری و دو نویسه‌ی پس از آن \\
\addlinespace
\cmd{dd} & حذف کامل سطر جاری \\
\addlinespace
\cmd{5dd} & حذف سطر جاری و چهار سطر بعدی \\
\addlinespace
\cmd{dW} & حذف از موقعیت فعلی مکان‌نما تا ابتدای کلمه‌ی بعدی \\
\addlinespace
\cmd{d\$} & حذف از موقعیت فعلی مکان‌نما تا انتهای سطر جاری \\
\addlinespace
\cmd{d0} & حذف از موقعیت فعلی مکان‌نما تا ابتدای مطلق سطر \\
\addlinespace
\cmd{d\^{}} & حذف از موقعیت فعلی مکان‌نما تا نخستین نویسه‌ی غیرفضای‌خالی سطر \\
\addlinespace
\cmd{dG} & حذف از سطر جاری تا انتهای فایل \\
\addlinespace
\cmd{d20G} & حذف از سطر جاری تا سطر بیستم فایل \\
\bottomrule
\end{tabularx}
\end{table}

برای تمرین، مکان‌نما را بر روی واژه‌ی \en{``It''} در سطر نخست قرار دهید. کلید \cmd{x} را چندین مرتبه بفشارید تا باقی‌مانده‌ی جمله حذف شود. سپس با فشردن متوالی کلید \cmd{u}، عملیات حذف را واگرد (\en{Undo}) نمایید تا متن به حالت اولیه بازگردد.

\begin{note}
\textbf{نکته:} در نسخه‌ی کلاسیک \cmd{vi} تنها یک سطح برای عملیات واگرد (\en{Undo}) پیش‌بینی شده است، در حالی که ویرایشگر بهبودیافته‌ی \cmd{vim} از سطوح متعدد \en{Undo} پشتیبانی می‌کند.
\end{note}

بیایید فرآیند حذف را این‌بار با استفاده از دستور \cmd{d} (مخفف \en{Delete}) بررسی کنیم. مکان‌نما را روی واژه‌ی \en{``It''} قرار داده و فرمان \cmd{dW} را صادر کنید تا این کلمه به‌طور کامل حذف شود:

\begin{latin}
\begin{lstlisting}
The quick brown fox jumps over the lazy dog. was cool.
Line 2
Line 3
Line 4
Line 5
\end{lstlisting}
\end{latin}

اکنون دستور \cmd{d\$} را تایپ نمایید. این دستور تمامی محتوا را از موقعیت کنونی مکان‌نما تا انتهای سطر جاری پاکسازی می‌کند:

\begin{latin}
\begin{lstlisting}
The quick brown fox jumps over the lazy dog.
Line 2
Line 3
Line 4
Line 5
\end{lstlisting}
\end{latin}

در مرحله بعد، دستور \cmd{dG} را وارد کنید. این فرمانِ قدرتمند، تمامی سطور را از محل استقرار مکان‌نما تا پایان فایل حذف خواهد کرد:

\begin{latin}
\begin{lstlisting}
~
~
~
~
~
\end{lstlisting}
\end{latin}

در نهایت، با سه بار فشردن کلید \cmd{u} (واگرد)، تمامی عملیات‌های حذف انجام‌شده را به ترتیب معکوس بازگردانی کنید تا فایل به وضعیت اولیه خود بازگردد.

\subsection{برش، کپی و درج متن}

در ویرایشگر \cmd{vi} دستور \cmd{d} علاوه بر حذف، عمل «برش» (\en{Cut}) را نیز انجام می‌دهد. هر بار که محتوایی را با این دستور حذف می‌کنید، متن مذکور در یک بافر موقت (مشابه به حافظه‌ی \en{Clipboard}) ذخیره می‌شود. سپس می‌توانید با استفاده از دستور \cmd{p} (حرف کوچک)، محتوای بافر را بعد از موقعیت مکان‌نما، و با دستور \cmd{P} (حرف بزرگ)، آن را قبل از موقعیت مکان‌نما درج (\en{Paste}) کنید.

دستور \cmd{y} (مخفف \en{Yank})، وظیفه‌ی کپی کردن متن را بر عهده دارد. منطق عملکرد این دستور کاملاً مشابه فرمان \cmd{d} است، با این تفاوت که متن در جای خود باقی می‌ماند و صرفاً نسخه‌ای از آن به بافر منتقل می‌شود. در جدول زیر، نمونه‌هایی از ترکیب دستور \cmd{y} با دستورات حرکتی (\en{Motions}) آمده است:

\begin{table}[htbp]
\centering
\caption{دستورات کپی متن (\en{Yank})}
\label{tab:vi-yank}
\begin{tabularx}{\textwidth}{l X}
\toprule
\textbf{دستور} & \textbf{دامنه‌ی کپی (\en{Scope})} \\
\midrule
\cmd{yy} & کپی کامل سطر جاری \\
\addlinespace
\cmd{5yy} & کپی سطر جاری و چهار سطر پس از آن \\
\addlinespace
\cmd{yW} & کپی از محل مکان‌نما تا ابتدای کلمه‌ی بعدی \\
\addlinespace
\cmd{y\$} & کپی از محل مکان‌نما تا انتهای سطر جاری \\
\addlinespace
\cmd{y0} & کپی از محل مکان‌نما تا ابتدای مطلق سطر \\
\addlinespace
\cmd{y\^{}} & کپی از محل مکان‌نما تا نخستین نویسه‌ی غیرفضای‌خالی سطر \\
\addlinespace
\cmd{yG} & کپی از سطر جاری تا انتهای فایل \\
\addlinespace
\cmd{y20G} & کپی از سطر جاری تا سطر بیستم فایل \\
\bottomrule
\end{tabularx}
\end{table}

بیایید فرآیند کپی و درج متن (\en{Copy \& Paste}) را عملاً بررسی کنیم. ابتدا مکان‌نما را بر روی سطر اول قرار داده و فرمان \cmd{yy} را برای کپی کردن کامل آن سطر صادر کنید. در مرحله بعد، با فشردن کلید \cmd{G} به آخرین سطر فایل بروید و کلید \cmd{p} را بفشارید تا محتوای کپی‌شده در سطری جدید، پایین‌تر از موقعیت فعلی درج شود:

\begin{latin}
\begin{lstlisting}
The quick brown fox jumps over the lazy dog. It was cool.
Line 2
Line 3
Line 4
Line 5
The quick brown fox jumps over the lazy dog. It was cool.
\end{lstlisting}
\end{latin}

همان‌طور که پیش‌تر اشاره شد، با فشردن کلید \cmd{u} می‌توانید این تغییر را واگرد کنید. اکنون در حالی که مکان‌نما همچنان روی سطر آخر قرار دارد، کلید \cmd{P} (حرف بزرگ) را بفشارید تا محتوا در سطری جدید، اما این‌بار بالاتر از موقعیت فعلی درج شود:

\begin{latin}
\begin{lstlisting}
The quick brown fox jumps over the lazy dog. It was cool.
Line 2
Line 3
Line 4
The quick brown fox jumps over the lazy dog. It was cool.
Line 5
\end{lstlisting}
\end{latin}

توصیه می‌شود سایر دستورات حرکتی و کپی مندرج در جدول فوق را نیز آزمایش کنید تا به تفاوت عملکرد \cmd{p} و \cmd{P} در سناریوهای مختلف مسلط شوید. در پایان تمرین، با استفاده از کلید \cmd{u} فایل را به وضعیت اولیه خود بازگردانید.

\subsection{ادغام سطور}

ویرایشگر \cmd{vi} دیدگاه بسیار سخت‌گیرانه‌ای نسبت به مفهوم «سطر» دارد. در این محیط، شما نمی‌توانید صرفاً با بردن مکان‌نما به انتهای سطر و حذف نویسه‌ی پنهانِ خط جدید (\en{Newline Character})، سطر فعلی را با سطر بعدی ادغام کنید. به همین منظور، \cmd{vi} دستور اختصاصی \cmd{J} (حرف بزرگ) را برای پیوستن سطور پیش‌بینی کرده است (دقت کنید که این دستور با کلید \cmd{j} که برای جابه‌جایی به سمت پایین است، اشتباه گرفته نشود).

چنانچه مکان‌نما را بر روی عبارت «\en{Line 3}» قرار داده و کلید \cmd{J} را بفشارید، سطر پایین به سطر فعلی الصاق شده و خروجی به صورت زیر تغییر می‌یابد:

\begin{latin}
\begin{lstlisting}
The quick brown fox jumps over the lazy dog. It was cool.
Line 2
Line 3 Line 4
Line 5
\end{lstlisting}
\end{latin}

همان‌طور که مشاهده می‌کنید، \cmd{vi} به‌صورت خودکار یک فاصله‌گذاری (\en{Space}) میان دو بخش ادغام‌شده ایجاد می‌کند تا یکپارچگی متن حفظ شود.

\section{جست‌وجو و جایگزینی}

ویرایشگر \cmd{vi} قابلیت‌های پیشرفته‌ای برای پیمایش مکان‌نما بر اساس الگوهای جست‌وجو در سطح یک سطر یا کل فایل فراهم می‌کند. علاوه بر این، امکان جایگزینی خودکار متن، با نظارت مستقیم کاربر یا بدون آن، میسر است.

\subsection{جست‌وجو در محدوده سطر}

دستور \cmd{f} (مخفف \en{find}) جهت جست‌وجوی نویسه‌ای خاص در سطر جاری به کار می‌رود و مکان‌نما را مستقیماً به نخستین رخداد آن نویسه منتقل می‌کند. به عنوان مثال، با اجرای دستور \cmd{fa}، مکان‌نما به اولین حرف \cmd{a} در سطر فعلی پرش می‌کند.

پس از اجرای جست‌وجوی اولیه، برای تکرار عملیات و انتقال به رخداد بعدی همان نویسه در همان سطر، کافی است کلید \cmd{;} (سمی‌کالن) را بفشارید.

\subsection{جست‌وجوی سراسری}

برای انتقال مکان‌نما به رخداد بعدی یک رشتهٔ متنی یا الگو (\en{Pattern}) در کل فایل، از دستور \cmd{/} استفاده می‌شود. منطق عملکرد این دستور دقیقاً مشابه فرآیندی است که پیش‌تر در ابزار \cmd{less} مشاهده کردیم. با فشردن کلید \cmd{/}، این نویسه در خط فرمان (پایین صفحه) ظاهر می‌شود؛ سپس عبارت مورد نظر را تایپ کرده و کلید \en{Enter} را بفشارید. پس از آن، مکان‌نما به نخستین انطباقِ پس از موقعیت فعلی پرش می‌کند. برای تکرار جست‌وجو و یافتن موارد بعدی، از دستور \cmd{n} استفاده می‌شود.

مثال زیر را در نظر بگیرید:

\begin{latin}
\begin{lstlisting}
The quick brown fox jumps over the lazy dog. It was cool.
Line 2
Line 3
Line 4
Line 5
\end{lstlisting}
\end{latin}

ابتدا مکان‌نما را در سطر نخست فایل قرار دهید. سپس دستور زیر را وارد کرده و \en{Enter} بزنید:

\begin{latin}
\begin{lstlisting}
/Line
\end{lstlisting}
\end{latin}

مکان‌نما بلافاصله به سطر دوم منتقل می‌شود. با فشردن کلید \cmd{n}، مکان‌نما به سطر سوم پرش خواهد کرد. با هر بار تکرار کلید \cmd{n}، پیمایش به سمت پایین ادامه می‌یابد تا زمانی که دیگر تطابقی یافت نشود.

تا این بخش، صرفاً از رشته‌های متنی ساده برای جست‌وجو استفاده کرده‌ایم؛ با این حال، ویرایشگر \cmd{vi} از عبارات باقاعده (\en{Regular Expressions}) نیز پشتیبانی می‌کند. این قابلیت، روشی بسیار منعطف و قدرتمند برای تعریف الگوهای متنی پیچیده است که در فصل ۱۹ به‌طور مفصل به بررسی آن خواهیم پرداخت.

\subsection{جست‌وجو و جایگزینی سراسری}

ویرایشگر \cmd{vi} از «حالت خط فرمان» (\en{Command-line Mode}) برای اجرای عملیات جایگزینی (\en{Substitution}) در محدوده‌ای مشخص از سطور یا کل فایل بهره می‌برد. برای مثال، جهت تغییر تمامی واژگان \en{``Line''} به \en{``line''} در سراسر سند، دستور زیر را وارد می‌کنیم:

\begin{latin}
\begin{lstlisting}
:%s/Line/line/g
\end{lstlisting}
\end{latin}

تحلیل اجزای این دستور در جدول زیر ارائه شده است:

\begin{table}[htbp]
\centering
\caption{نحو (\en{Syntax}) جست‌وجو و جایگزینی}
\label{tab:vi-substitute}
\begin{tabularx}{\textwidth}{l X}
\toprule
\textbf{مؤلفه} & \textbf{مفهوم فنی} \\
\midrule
\cmd{:} & ورود به حالت خط فرمان (\en{Command-line Mode}) \\
\addlinespace
\cmd{\%} & تعیین محدوده‌ی عملیاتی؛ کاراکتر \cmd{\%} نماد اختصاری «از سطر اول تا سطر آخر» است. مقادیر معادل آن می‌تواند \cmd{1,5} (برای یک فایل ۵ سطری) یا \cmd{1,\$} (از سطر اول تا انتها) باشد. در صورت عدم درج محدوده، عملیات فقط روی سطر جاری اعمال می‌شود. \\
\addlinespace
\cmd{s} & نشان‌دهنده‌ی نوع عملیات؛ جایگزینی (مخفف \en{Substitute}). \\
\addlinespace
\cmd{/Line/line/} & تعیین الگوی هدف (جست‌وجو) و عبارت جایگزین. \\
\addlinespace
\cmd{g} & مخفف \en{Global}؛ تضمین می‌کند که تمامی تکرارهای الگوی مورد نظر در هر سطر جایگزین شوند. در صورت حذف این پارامتر، تنها اولین رخداد در هر سطر تغییر می‌یابد. \\
\bottomrule
\end{tabularx}
\end{table}

پس از اجرای دستور فوق، محتوای فایل به صورت زیر تغییر می‌یابد:

\begin{latin}
\begin{lstlisting}
The quick brown fox jumps over the lazy dog. It was cool.
line 2
line 3
line 4
line 5
\end{lstlisting}
\end{latin}

همچنین می‌توان فرآیند جایگزینی را به‌گونه‌ای پیکربندی کرد که برای هر مورد از کاربر تأییدیه دریافت کند. این عمل با افزودن پارامتر \cmd{c} (مخفف \en{Confirmation}) به انتهای دستور میسر می‌شود. به عنوان مثال:

\begin{latin}
\begin{lstlisting}
:%s/line/Line/gc
\end{lstlisting}
\end{latin}

این دستور فایل را به وضعیت پیشین بازمی‌گرداند، اما پیش از هر تغییر، \cmd{vi} متوقف شده و پیامی مشابه زیر را برای دریافت تاییدیه نمایش می‌دهد:

\begin{latin}
\begin{lstlisting}
replace with Line (y/n/a/q/l/^E/^Y)?
\end{lstlisting}
\end{latin}

عملکرد هر یک از این گزینه‌ها در جدول زیر تشریح شده است:

\begin{table}[htbp]
\centering
\caption{کلیدهای کنترل تأییدیه در جایگزینی}
\label{tab:vi-confirm}
\begin{tabularx}{\textwidth}{l X}
\toprule
\textbf{کلید} & \textbf{عملکرد فنی} \\
\midrule
\cmd{y} & تایید و انجام جایگزینی (\en{Yes}) \\
\addlinespace
\cmd{n} & رد جایگزینی برای مورد فعلی و پرش به مورد بعد (\en{No}) \\
\addlinespace
\cmd{a} & پذیرش جایگزینی برای مورد فعلی و تمامی موارد باقی‌مانده (\en{All}) \\
\addlinespace
\cmd{q} یا \en{Esc} & توقف آنی و خروج از فرآیند جایگزینی (\en{Quit}) \\
\addlinespace
\cmd{l} & انجام جایگزینی فعلی و سپس خروج از فرآیند (\en{Last}) \\
\addlinespace
\cmd{Ctrl+e} / \cmd{Ctrl+y} & پیمایش به پایین و بالا جهت مشاهدهٔ بافت گسترده‌تر متن و اتخاذ تصمیم دقیق‌تر \\
\bottomrule
\end{tabularx}
\end{table}

\section{ویرایش هم‌زمان چندین فایل}

در بسیاری از موارد، ویرایش هم‌زمان چند فایل مزیتی راهبردی محسوب می‌شود؛ به‌ویژه زمانی که نیاز به اعمال تغییرات هم‌زمان در چندین فایل داشته باشید یا بخواهید قطعاتی از متن را میان فایل‌های مختلف جابه‌جا کنید. در ویرایشگر \cmd{vi}، باز کردن چندین فایل به‌طور هم‌زمان بسیار ساده است و تنها کافی است نام آن‌ها را به‌صورت متوالی در خط فرمان درج کنید:

\begin{latin}
\begin{lstlisting}
vi file1 file2 file3...
\end{lstlisting}
\end{latin}

برای شروع این تمرین، ابتدا باید از نشست (\en{session}) فعلی \cmd{vi} خارج شوید. جهت ذخیره‌سازی تغییرات و خروج قطعی، دستور زیر را وارد نمایید:

\begin{latin}
\begin{lstlisting}
:wq
\end{lstlisting}
\end{latin}

سپس، برای ایجاد یک فایل کمکی در دایرکتوری خانگی (\en{Home Directory})، خروجی دستور \cmd{ls} را در فایلی به نام \file{ls-output.txt} بازنویسی می‌کنیم:

\begin{latin}
\begin{lstlisting}
[me@linuxbox ~]$ ls -l /usr/bin > ls-output.txt
\end{lstlisting}
\end{latin}

اکنون فایل پیشین و فایل جدید را به‌طور هم‌زمان با \cmd{vi} فراخوانی می‌کنیم:

\begin{latin}
\begin{lstlisting}
[me@linuxbox ~]$ vi foo.txt ls-output.txt
\end{lstlisting}
\end{latin}

پس از اجرا، \cmd{vi} کار خود را آغاز کرده و نخستین فایل فهرست‌شده، یعنی \file{foo.txt}، در محیط ویرایشگر بارگذاری می‌شود:

\begin{latin}
\begin{lstlisting}
The quick brown fox jumps over the lazy dog. It was cool.
Line 2
Line 3
Line 4
Line 5
\end{lstlisting}
\end{latin}

\subsection{مدیریت بافرها و پیمایش فایل‌ها}

برای جابه‌جایی به فایل بعدی در فهرست فایل‌های باز، از دستور \cmd{ex} زیر استفاده می‌شود:

\begin{latin}
\begin{lstlisting}
:bn
\end{lstlisting}
\end{latin}

به همین ترتیب، جهت بازگشت به فایل پیشین، دستور زیر را به کار می‌بریم:

\begin{latin}
\begin{lstlisting}
:bp
\end{lstlisting}
\end{latin}

ویرایشگر \cmd{vi} از رفتاری حفاظتی پیروی می‌کند که طبق آن، در صورت وجود تغییرات ذخیره‌نشده در فایل جاری، اجازه جابه‌جایی به فایل دیگر را صادر نمی‌کند. برای اجبار \cmd{vi} به تغییر فایل و نادیده گرفتن تغییرات (بدون ذخیره‌سازی)، باید کاراکتر تعجب \cmd{!} را به انتهای دستور اضافه نمایید (مانند \cmd{:bn!}).

علاوه بر روش‌های فوق، \cmd{vim} (و نسخه‌های پیشرفته \cmd{vi}) قابلیت‌های قدرتمندتری را برای مدیریت هم‌زمان فایل‌ها در قالب «بافرها» فراهم کرده‌اند. برای مشاهده فهرست تمامی فایل‌های بارگذاری شده، دستور زیر را وارد کنید:

\begin{latin}
\begin{lstlisting}
:buffers
\end{lstlisting}
\end{latin}

پس از اجرا، فهرستی مشابه نمونه زیر در پایین صفحه ظاهر می‌شود:

\begin{latin}
\begin{lstlisting}
:buffers
  1 %a  "foo.txt"                     line 1
  2     "ls-output.txt"               line 0
Press ENTER or type command to continue
\end{lstlisting}
\end{latin}

برای انتقال مستقیم به یک بافر خاص، دستور \cmd{:buffer} را به همراه شماره اختصاصی آن بافر وارد کنید. برای مثال، جهت جابه‌جایی از بافر ۱ به بافر ۲، دستور زیر را تایپ نمایید:

\begin{latin}
\begin{lstlisting}
:buffer 2
\end{lstlisting}
\end{latin}

در این وضعیت، محتوای فایل دوم روی نمایشگر ظاهر خواهد شد. شایان ذکر است که دستورات \cmd{:bn} (مخفف \en{buffer next}) و \cmd{:bp} (مخفف \en{buffer previous}) نیز همچنان برای پیمایش ترتیبی میان بافرها قابل استفاده هستند.

\subsection{فراخوانی فایل‌های جدید جهت ویرایش}

علاوه بر باز کردن هم‌زمان فایل‌ها هنگام شروع به کار، این امکان وجود دارد که در حین نشست ویرایش فعلی، فایل‌های جدیدی را نیز به محیط اضافه کنید. دستور \cmd{:e} (مخفف \en{edit}) در حالت خط فرمان (\en{Command-line Mode})، به همراه نام فایل مورد نظر، فایل جدیدی را بارگذاری می‌کند. برای تمرین این قابلیت، ابتدا از نشست فعلی خارج شده و به خط فرمان بازگردید.

سپس \cmd{vi} را تنها با فراخوانی یک فایل اجرا کنید:

\begin{latin}
\begin{lstlisting}
[me@linuxbox ~]$ vi foo.txt
\end{lstlisting}
\end{latin}

برای افزودن فایل دوم به محیط ویرایش، دستور زیر را صادر کنید:

\begin{latin}
\begin{lstlisting}
:e ls-output.txt
\end{lstlisting}
\end{latin}

در این مرحله، فایل جدید بر روی صفحه ظاهر می‌شود. توجه داشته باشید که فایل نخست همچنان در حافظه‌ی موقت (\en{Buffer}) باقی مانده است؛ برای اطمینان از این موضوع، می‌توانید فهرست بافرهای فعال را بازبینی کنید:

\begin{latin}
\begin{lstlisting}
:buffers
  1 #   "foo.txt"                     line 1
  2 %a  "ls-output.txt"               line 0
Press ENTER or type command to continue
\end{lstlisting}
\end{latin}

\begin{note}
\textbf{نکته:} در خروجی دستور فوق، نشانگر \cmd{\%a} نشان‌دهنده‌ی بافر فعال (\en{Active}) و نشانگر \cmd{\#} نشان‌دهنده‌ی بافر قبلی است که همچنان در حافظه بارگذاری شده است.
\end{note}

\subsection{انتقال محتوا میان فایل‌ها}

هنگام کار با چندین فایل، اغلب نیاز داریم بخشی از محتوای یک سند را به سند دیگری منتقل کنیم. این فرآیند با استفاده از دستورات \en{Yank} (کپی) و \en{Put} (درج) که پیش‌تر آموختیم، به‌سادگی امکان‌پذیر است.

برای مشاهده‌ی این عملکرد، در حالی که هر دو فایل باز هستند، ابتدا با دستور زیر به بافر نخست (فایل \file{foo.txt}) بازگردید:

\begin{latin}
\begin{lstlisting}
:buffer 1
\end{lstlisting}
\end{latin}

در این حالت، محتوای زیر روی صفحه نمایان می‌شود:

\begin{latin}
\begin{lstlisting}
The quick brown fox jumps over the lazy dog. It was cool.
Line 2
Line 3
Line 4
Line 5
\end{lstlisting}
\end{latin}

اکنون مکان‌نما را روی سطر اول قرار داده و با فشردن کلید \cmd{yy}، آن سطر را در بافر کپی (\en{Yank}) کنید. سپس با استفاده از دستور زیر به بافر دوم جابه‌جا شوید:

\begin{latin}
\begin{lstlisting}
:buffer 2
\end{lstlisting}
\end{latin}

در این مرحله، صفحه‌ی نمایش فهرستی از فایل‌ها را (مشابه خروجی زیر) نشان می‌دهد:

\begin{latin}
\begin{lstlisting}
total 343700
-rwxr-xr-x 1 root root       31316 2017-12-05 08:58 [
-rwxr-xr-x 1 root root        8240 2017-12-09 13:39 411toppm
-rwxr-xr-x 1 root root      111276 2018-01-31 13:36 a2p
-rwxr-xr-x 1 root root       25368 2016-10-06 20:16 a52dec
-rwxr-xr-x 1 root root       11532 2017-05-04 17:43 aafire
-rwxr-xr-x 1 root root        7292 2017-05-04 17:43 aainfo
\end{lstlisting}
\end{latin}

مکان‌نما را به سطر اول انتقال دهید و با فشردن کلید \cmd{p}، سطری را که از فایل پیشین کپی کرده بودید، در اینجا درج (\en{Paste}) کنید:

\begin{latin}
\begin{lstlisting}
total 343700
The quick brown fox jumps over the lazy dog. It was cool.
-rwxr-xr-x 1 root root       31316 2017-12-05 08:58 [
-rwxr-xr-x 1 root root        8240 2017-12-09 13:39 411toppm
-rwxr-xr-x 1 root root      111276 2018-01-31 13:36 a2p
-rwxr-xr-x 1 root root       25368 2016-10-06 20:16 a52dec
-rwxr-xr-x 1 root root       11532 2017-05-04 17:43 aafire
-rwxr-xr-x 1 root root        7292 2017-05-04 17:43 aainfo
\end{lstlisting}
\end{latin}

\subsection{درج محتوای یک فایل در فایل دیگر}

علاوه بر کپی قطعات متنی، ویرایشگر \cmd{vi} این امکان را فراهم می‌کند تا محتوای کامل یک فایل خارجی را مستقیماً درون فایلی که در حال ویرایش آن هستید، درج کنید. برای مشاهده عملی این قابلیت، ابتدا از نشست جاری خارج شده و نشست جدیدی را تنها با یک فایل آغاز کنید:

\begin{latin}
\begin{lstlisting}
[me@linuxbox ~]$ vi ls-output.txt
\end{lstlisting}
\end{latin}

مجدداً فهرست فایل‌ها را مشاهده خواهید کرد:

\begin{latin}
\begin{lstlisting}
total 343700
-rwxr-xr-x 1 root root       31316 2017-12-05 08:58 [
-rwxr-xr-x 1 root root        8240 2017-12-09 13:39 411toppm
-rwxr-xr-x 1 root root      111276 2018-01-31 13:36 a2p
-rwxr-xr-x 1 root root       25368 2016-10-06 20:16 a52dec
-rwxr-xr-x 1 root root       11532 2017-05-04 17:43 aafire
-rwxr-xr-x 1 root root        7292 2017-05-04 17:43 aainfo
\end{lstlisting}
\end{latin}

مکان‌نما را به سطر سوم انتقال دهید و دستور زیر را در حالت خط فرمان (\en{Command-line Mode}) وارد کنید:

\begin{latin}
\begin{lstlisting}
:r foo.txt
\end{lstlisting}
\end{latin}

دستور \cmd{:r} (مخفف \en{read})، محتوای فایل مشخص‌شده را دقیقاً در سطر زیرینِ موقعیت فعلی مکان‌نما درج می‌کند. اکنون چیدمان متن در صفحه به صورت زیر تغییر می‌یابد:

\begin{latin}
\begin{lstlisting}
total 343700
-rwxr-xr-x 1 root root       31316 2017-12-05 08:58 [
-rwxr-xr-x 1 root root        8240 2017-12-09 13:39 411toppm
The quick brown fox jumps over the lazy dog. It was cool.
Line 2
Line 3
Line 4
Line 5
-rwxr-xr-x 1 root root      111276 2018-01-31 13:36 a2p
-rwxr-xr-x 1 root root       25368 2016-10-06 20:16 a52dec
-rwxr-xr-x 1 root root       11532 2017-05-04 17:43 aafire
-rwxr-xr-x 1 root root        7292 2017-05-04 17:43 aainfo
\end{lstlisting}
\end{latin}

\section{روش‌های ذخیره‌سازی و خروج}

مانند سایر قابلیت‌های \cmd{vi}، برای ذخیره‌سازی فایل‌های ویرایش‌شده نیز چندین روش عملیاتی وجود دارد. پیش‌تر با دستور \cmd{:w} آشنا شدیم، اما دستورات کاربردی دیگری نیز در این زمینه در دسترس هستند:

در حالت عادی (\en{Normal Mode})، با فشردن کلیدهای ترکیبی \cmd{ZZ}، فایل جاری ذخیره شده و ویرایشگر بسته می‌شود. به همین ترتیب، فرمان ترکیبی \cmd{:wq} با ادغام دو دستور \cmd{:w} (نوشتن/ذخیره) و \cmd{:q} (خروج)، فرآیند ذخیره‌سازی و خروج را به‌صورت هم‌زمان به انجام می‌رساند.

دستور \cmd{:w} همچنین قابلیت پذیرش یک نام فایل اختیاری را دارد که عملکردی مشابه گزینه \en{``Save As\ldots''} در سایر ویرایشگرها ارائه می‌دهد. برای نمونه، اگر در حال ویرایش فایل \file{foo.txt} هستید و قصد دارید یک نسخه کپی از آن را با نام \file{foo1.txt} ذخیره کنید، دستور زیر را به کار می‌برید:

\begin{latin}
\begin{lstlisting}
:w foo1.txt
\end{lstlisting}
\end{latin}

\begin{note}
\textbf{نکته:} توجه داشته باشید که این عمل صرفاً محتوا را در فایلی با نام جدید بازنویسی می‌کند، اما بافر جاری ویرایشگر را به فایل جدید منتقل نمی‌کند. بنابراین، شما همچنان در حال ویرایش فایل اصلی (\file{foo.txt}) باقی خواهید ماند و نه نسخه جدید (\file{foo1.txt}).
\end{note}

\section{پشتیبانی پوسته Bash از سبک ویرایش vi}

در فصل هشتم، روش‌های مختلف ویرایش در خط فرمان را بررسی کردیم. به‌صورت پیش‌فرض، پوسته \cmd{bash} از کلیدهای میان‌بر الهام‌گرفته از ویرایشگر \cmd{emacs} بهره می‌برد؛ با این حال، \cmd{bash} امکان سوییچ به حالت ویرایش منطبق بر سبک \cmd{vi} را نیز فراهم کرده است. این قابلیت با اجرای دستور زیر فعال می‌شود:

\begin{latin}
\begin{lstlisting}
[me@linuxbox ~]$ set -o vi
\end{lstlisting}
\end{latin}

پس از اجرای این دستور، می‌توانید از بسیاری از فرامین ویرایشی \cmd{vi} که تاکنون آموخته‌اید، در محیط ترمینال استفاده کنید. برای آزمایش، عبارت زیر را در خط فرمان تایپ کنید:

\begin{latin}
\begin{lstlisting}
[me@linuxbox ~]$ the quick brown fox jumps over the lazy dog
\end{lstlisting}
\end{latin}

در این وضعیت، کماکان امکان جابه‌جایی با کلیدهای جهت‌نما و تایپ نویسه‌ها فراهم است؛ چرا که \cmd{bash} هر خط فرمان جدید را در حالت درج (\en{Insert Mode}) آغاز می‌کند (این رفتار با \cmd{vim} که به‌طور پیش‌فرض در حالت عادی (\en{Normal Mode}) شروع به کار می‌کند، تفاوت دارد).

برای بهره‌گیری از قابلیت‌های پیشرفته، باید به حالت عادی (\en{Normal Mode}) وارد شوید. با فشردن کلید \en{ESC}، از حالت درج خارج شده و تمامی دستورات جابه‌جایی مکان‌نما، کپی (\en{Yank})، حذف (\en{Delete}) و درج (\en{Paste}) دقیقاً مانند ویرایش یک فایل تک‌خطی در \cmd{vi} عمل خواهند کرد. جهت بازگشت به حالت ویرایش متنی، از دستورات استاندارد حالت عادی مانند \cmd{i} یا \cmd{A} استفاده کنید.

پیکربندی \cmd{bash} بر روی حالت \cmd{vi}، راهکاری هوشمندانه برای تقویت حافظه‌ی عضلانی و مهارت‌های صفحه‌کلیدی شما در این ویرایشگر است و نیاز به حفظ کردن دو مجموعه متفاوت از کلیدهای میان‌بر را از بین می‌برد. برای دائمی کردن این تنظیم، می‌توانید دستور \cmd{set -o vi} را به فایل پیکربندی \file{.bashrc} خود اضافه کنید.

چنانچه تمایل به بازگشت به حالت پیش‌فرض (\cmd{emacs}) داشتید، دستور زیر را اجرا کنید:

\begin{latin}
\begin{lstlisting}
[me@linuxbox ~]$ set -o emacs
\end{lstlisting}
\end{latin}

\section{جمع‌بندی}

با کسب این مهارت‌های کلیدی، اکنون شما قادر هستید بخش اعظم وظایف ویرایشی لازم برای مدیریت و نگهداری یک سیستم لینوکسی را به انجام برسانید. ممارست در استفاده از ویرایشگر \cmd{vim}، در بلندمدت سرمایه‌گذاری ارزشمندی خواهد بود؛ چرا که ویرایشگرهای خانواده \cmd{vi} عمیقاً در فرهنگ و فلسفه یونیکس ریشه دوانده‌اند. ردپای طراحی منحصربه‌فرد این ویرایشگر در ابزارهای متعددی مشاهده می‌شود که برنامه‌ی \cmd{less} تنها نمونه‌ای از این تأثیرپذیری گسترده است.

\section{منابع مطالعه تکمیلی}

مطالب ارائه‌شده در این فصل، تنها مقدمه‌ای بر قابلیت‌های گسترده ویرایشگر \cmd{vim} است. برای افزایش تسلط خود، مطالعه منابع زیر پیشنهاد می‌شود:

کتاب \en{``Vim, with Vigor''} – دنباله‌ای بر این آموزش در وب‌سایت \en{LinuxCommand.org} که دانش خواننده را به سطح متوسط ارتقاء می‌دهد:

\begin{latin}
\url{https://linuxcommand.org/lc3_adv_vimvigor.php}
\end{latin}

کتاب \en{``Learning The vi Editor''} – یک ویکی‌کتاب از ویکی‌پدیا که راهنمایی خلاصه برای \cmd{vi} و ویرایشگرهای مشابه آن نظیر \cmd{vim} ارائه می‌کند:

\begin{latin}
\url{http://en.wikibooks.org/wiki/Vi}
\end{latin}

مقاله ویکی‌پدیا درباره «بیل جوی» (\en{Bill Joy})، خالق ویرایشگر \cmd{vi}:

\begin{latin}
\url{http://en.wikipedia.org/wiki/Bill_Joy}
\end{latin}

مقاله ویکی‌پدیا درباره «برام مولنار» (\en{Bram Moolenaar})، نویسنده ویرایشگر \cmd{vim}:

\begin{latin}
\url{http://en.wikipedia.org/wiki/Bram_Moolenaar}
\end{latin}